\section{Independent orbit proofs and reusable marked-union queries}
\label{sec:orbit-certificates}

Report 47 supplies an independent verifier for the six AHT reduction rules
and a reduction from marked-component incidence to ordinary orbit counts.
We reconcile these with the sharper kernel of
Section~\ref{sec:native-interval-orbits}, instead of installing a second
production orbit scheduler. The maintained extension records optional local
proofs, replays both merger conventions, computes complete incidence
histograms, and shares repeated marked-union queries. These are local
components of a possible compressed hierarchy, not a recognition procedure.

\subsection{What the verifier establishes}

\path{../fast/fastunknot/interval_orbit_verify.py} imports neither the orbit
producer nor its scheduling helpers. It checks the universe size and the
ordered, canonical input pairing list, then applies six independently
implemented local rules. A certificate must finish with an empty universe
and no pairings. Its claimed count must equal the number of singleton
orbits removed by static contractions. A truncated valid prefix proves no
count. Boolean values are rejected in integer fields, and exact hexadecimal
transport is accepted without changing Python's global decimal limit.

Identity deletion removes only reflexive relations. Static contraction
recomputes the entire complement of all pairing domains and ranges,
compares the supplied gaps, and shifts remaining endpoints by their exact
deleted rank. Reflection trimming retains one arrow for each two-point
orbit and leaves a fixed midpoint available for subsequent counting.
A periodic merger checks both periods and the union hull. Version one
requires overlap $p+q$; version two requires $p+q-\gcd(p,q)$, using the
proved criterion in the preceding section. Changing a sharp-only trace
to version one therefore makes it fail replay.

Transmission records the powers applied to the source and image intervals
of a target pairing. The verifier checks that these partial inverses are
defined throughout each interval. For a positive translation of distance
$p$, power $s$ is legal only when the original interval is in the range
and its last intermediate left endpoint remains in that range. Monotonicity
then covers every intermediate power with one inequality; the cost depends
on the bit length of $s$, not on $s$. A trimmed reflection permits one
inverse step. The verifier composes the two affine maps independently and
normalizes their resulting intervals.

Suffix truncation requires the carrier to end at the universe boundary,
the removed suffix to lie in its range, and every other pairing to avoid
that suffix. A positive carrier has strictly positive displacement, so
iteration reaches a surviving earlier representative. A reflection must
already have disjoint domain and range. The verifier does not require the
producer's preferred scheduling choices: any sequence satisfying these
local conditions is sound. Scheduling is needed for the producer's
polynomial running-time theorem, not for the validity of a supplied trace.

Replay is polynomial in the explicit trace length, pairing-list size and
integer bit lengths. The producer's polynomial cycle bound gives a
polynomial-size trace: a cycle emits polynomially many deletion, gap, merge,
transmission and truncation records, each with polynomially many binary
integers. This is not a hostile-input memory cap; callers can limit
serialized size and use cooperative cancellation. Production recording is
optional. Incomplete runs expose neither a partial count nor a certificate.

\subsection{Coning a marked union}

Let $A$ be a nonempty union of half-open marked intervals in the orbit
universe. For each interval $[a,b)$ of length at least two, adjoin the
translation $[a,b-2]\to[a+1,b-1]$. Its points become one class. Join one
point of each interval to a common marked anchor by singleton pairings.
There are at most twice as many new pairings as marked intervals, with no
expansion into individual points.

If $C$ is the original orbit count and $C_A$ is the count after these
identifications, then the number of original orbits meeting $A$ is
\[
 H(A)=C-C_A+1.
\]
Indeed, exactly those original orbits meeting $A$ become one class; all
others retain their identities. The empty union has $H(\varnothing)=0$
and requires no coned query. This argument allows partial translations,
reflections, intersecting marks, static points and disconnected supports.

For ordered ports $A_1,\ldots,A_r$, let $h(T)$ count orbits meeting exactly
the ports indexed by $T\subseteq[r]$. Define
\[
 F(S)=H\!\left(\bigcup_{i\in S}A_i\right),\qquad
 G(U)=C-F([r]\setminus U).
\]
An orbit contributes to $G(U)$ precisely when its signature is a subset
of $U$, so $G(U)=\sum_{T\subseteq U}h(T)$. The producer recovers $h$ by
Boolean M\"obius inversion in $O(r2^r)$ integer subtractions. It retains
$h(\varnothing)$, the completely unmarked component count.

\subsection{Sharing queries and proofs without changing the output}

Different port subsets can describe the same marked set. The implementation
sorts intervals, merges overlap and adjacency, and uses the resulting
ordered disjoint union as its cache key. Equal unions introduce exactly
the same coning pairings in the same order. Consequently one result and
one proof can serve every mask with that union; neither an approximate
shape signature nor equality of counts is used to justify reuse.

If all $r$ nonempty ports coincide, there is one distinct nonempty union:
only two orbit queries are needed, including the original count. If the
ports form a strict nonempty nested chain, there are $r$ distinct unions
and $r+1$ queries. Disjoint nonempty ports can have $2^r-1$ distinct unions,
so no uniform query reduction follows. The output remains a dense array
of $2^r$ counts, even when query reuse is substantial.

The new incidence certificate has version two. It contains one baseline
orbit proof, one proof per distinct nonempty union, and an index array
binding every subset mask to its proof. Empty unions have no proof index.
The verifier reconstructs the marked unions from the supplied ports,
rejects rebinding one proof to a different union, checks every distinct
orbit proof, and rejects unused or duplicate proof assignments. To check
the histogram it computes its \emph{forward} subset sums, rather than
repeating the producer's inverse transform. This catches incorrect
histograms independently of inversion code. Input normalization and the
explicit coning construction are shared; their correctness is also checked
against literal marked graphs in the small-instance tests.

Let $M$ be the total number of marked intervals and $U$ the number of
distinct nonempty unions. If $P(k,B)$ bounds the orbit kernel for $k$
pairings and $B$-bit endpoints, a coarse bound is
\[
 (U+1)P(k+2M,B)
 +O\!\left(2^r M\log(M+1)B+r2^r B\right).
\]
The first term pays for the baseline and distinct coned queries; the second
pays for union construction, cache keys and dense inversion. For hash-based
lookups this uses ordinary expected dictionary costs; balanced ordered
keys would add polynomial comparison factors without changing the
$2^r\operatorname{poly}(k+M+r,B)$ bound. Certificate storage contains one
trace per distinct union and $O(2^r\log(U+2))$ bits of proof indices,
in addition to the dense counts and input records.

Polynomially bounded $k,M,B$ give a polynomial query for $r=O(\log n)$,
and a quasi-polynomial query for $r=O((\log n)^2)$. Neither condition has
been established for a complete unknot hierarchy. The weighted AHT
extension offers a different route to sparse component-weight information;
this dense reduction is not a new general complexity bound for weighted
orbits. It implements one useful aggregate query over the unweighted kernel.

Unlike the delivery's separate allowance for each coned query, the maintained
interface shares one \texttt{max\_cycles} budget across all orbit calls.
Exhaustion returns \texttt{INCONCLUSIVE} with work statistics and no histogram
or certificate. Port validation and Boolean inversion also poll the global
callback. The default explicit port cap is 12. Raising it permits the
corresponding exponential output cost; the cap is not a theorem about
geometric input. A signature records set incidence only. It omits point
multiplicity, cyclic order, slope and the action of future attachment maps.

\subsection{Validation and measured costs}

Nineteen new verifier tests adapt the report's adversarial replay suite and
add sharp-threshold version checks, a forged merger below the exact threshold,
1,000 randomized comparisons with literal graphs, and a transmission of
$2^{16000}-1$ steps stored in a single record. They compare traced and
untraced results and operation statistics, exercise exact JSON transport,
and disable producer functions while replaying proofs. Malformed inputs,
changed source bindings, invalid local operations and unfinished traces
are rejected. Seven incidence tests compare 250 random marked graphs and
18 actual normal-arc systems against literal component signatures, including
parallel meridians, overlapping and coincident ports, huge static universes,
shared budget boundaries and certificate mutations.
All 831 maintained tests pass with Regina in 179.802 seconds. A separate
cross-replay audit imports the unchanged report-47 producer and verifier
under another package namespace. For 1,000 interval systems, the maintained
verifier accepts the delivery's classical proofs and the delivery's verifier
accepts our classical proofs; our sharp proofs also pass independent replay.
The audit exercises all six rule types. One hundred marked systems agree
with the delivery's complete incidence histograms, with both proof formats
checked by their respective verifiers. Evidence and source hashes are in
\path{data/orbit-certificate-cross-replay.json}.

\begin{center}\small
\begin{tabular}{@{}lrrrrr@{}}
\toprule
Marked system & Fresh & Reuse & Certified & Ratio & Queries \\
\midrule
Static, coincident & 8.13 & 1.41 & 3.16 & 5.76 & 256/2 \\
Static, nested & 8.95 & 1.41 & 3.00 & 6.09 & 256/9 \\
Static, disjoint & 36.78 & 38.81 & 67.63 & 0.95 & 256/256 \\
Meridian, coincident & 128.47 & 3.79 & 5.76 & 33.31 & 64/2 \\
Meridian, nested & 52.57 & 10.06 & 15.31 & 5.26 & 64/7 \\
Meridian, disjoint & 61.32 & 61.59 & 98.67 & 0.99 & 64/64 \\
Parallel, coincident & 2410.80 & 40.47 & 68.65 & 61.82 & 64/2 \\
Parallel, nested & 1828.16 & 137.31 & 242.90 & 13.24 & 64/7 \\
Parallel, disjoint & 800.09 & 782.48 & 1357.14 & 1.02 & 64/64 \\
\bottomrule
\end{tabular}
\end{center}
Median milliseconds per complete incidence query. Certified includes proof construction and independent replay; ratio is the median paired fresh/reuse time. Queries are fresh/reuse.


The incidence benchmark uses the same maintained orbit kernel for both
the fresh-query ablation and the cached implementation, isolating reuse
from the difference between the reports' two orbit schedulers. There are
nine workloads, four shuffled arms, one warmup and five measured rounds:
36 warmup and 180 measured incidence calls. The fourth arm is an identical
cached control. Seed 261008483 fixes the order. Each complete histogram is
compared outside its timed region; the certified arm includes local proof
construction and independent verification in its timing. The inputs use
eight marks on a static $2^{4000}$-point universe, six marks on the normal
arc graph of an eight-tetrahedron meridian, or six marks on $2^{500}$
parallel copies of that meridian. Marks are chosen coordinate intervals,
not claimed outputs of a geometric attachment-port extractor.

Coincident marks reduce queries from 256 to two in the static case and
from 64 to two in the native cases. The native unscaled query improves
from 128.5 to 3.8 ms, with paired ratio 33.308. The large parallel-copy
query improves from 2.411 seconds to 40.5 ms, with paired ratio 61.817;
including proof generation and replay takes 68.6 ms. Ratios can exceed
the ratio of total query counts because the unchanged baseline count
is cheaper than a coned count. Nested marks also benefit, while disjoint
marks leave every union distinct. The latter show no reliable improvement
and can incur cache overhead. All timing arms and proof-size costs are
retained in \path{../fast/results/orbit_certificates_20261008.json};
for example, the disjoint parallel-copy certificate is about 13.3 MB
in the measured JSON transport. Binary input does not imply a tiny trace.

The same driver separately measures ordinary counts, optional recording,
and recording plus replay against the pre-certificate kernel pinned at
\texttt{8cd6dcfcd}. Its initial small-query samples had divergent identical
controls, so a separate batched follow-up retains those results and measures
fresh batches of the four workloads. It uses seed 261008485, seven measured
rounds, and 1,000 repeats for each tiny translation/reflection query, 30
for the 64-pair chain and ten for the native 32-tetrahedron arc query.
Its raw results are in \path{../fast/results/orbit_trace_cost_20261008.json}.
On the native arc query, the old/current paired ratio is 0.976: the optional
recording branches have a small default-path cost. Median times are 33.34 ms
before the change, 33.92 ms without recording, 37.48 ms with recording and
42.88 ms with recording plus replay. The translation and chain ratios are
0.968 and 0.993. Reflection samples retain some variability, so no precise
default-path gain is claimed there. Paired identical-current ratios in the
follow-up range from 0.999 to 1.026. Proof production is optional precisely
because its time and storage are real costs. These local costs do not alter
ordinary knot dispatch, which does not currently call the orbit interface.

The remaining normal-disc certificate adapter and sparse two-meridian seed
search in report 47 require separate review. The report's source snapshots
are unchanged. The maintained coning and replay code establish assertions
about supplied interval relations, including relations extracted by our
validated normal-arc adapter; they do not establish diagram-to-exterior
provenance or search completeness.
